\section{Inteligencia cognitiva: el siguiente paso después de la inteligencia perceptiva}

Entre 2015 y 2016, la generación anterior de empresas de IA se centraba sobre todo en llevar a la práctica la CV, el NLP temprano y el aprendizaje profundo. Zhang Peng (张鹏) cuenta que en ese momento se dieron cuenta de que la inteligencia perceptiva tenía un techo: una vez que el reconocimiento facial superó al ser humano, seguir mejorándolo tenía un sentido limitado, y su efecto en la industria también se acercaba a su límite. Por eso propusieron la «inteligencia cognitiva» como el siguiente escalón hacia la inteligencia artificial general.

\begin{figure}[H]
\centering
\includegraphics[width=0.94\textwidth]{cognitive-intelligence-ladder.png}
\caption{De la inteligencia perceptiva a la inteligencia cognitiva, y de ahí a los grandes modelos, los Agent y la AGI.}
\end{figure}

\begin{knowledgebox}{Lectura de la figura: la inteligencia cognitiva es una definición por etapas}
En la figura, la inteligencia cognitiva no es la AGI en sí, sino el siguiente paso después de la inteligencia perceptiva. Busca llevar la IA de las tareas de reconocimiento, clasificación y percepción a la comprensión, el razonamiento, el conocimiento, el lenguaje y las tareas complejas. Más tarde, los grandes modelos se convirtieron en la forma tecnológica principal de esa dirección.
\end{knowledgebox}

\subsection{Por qué no definir directamente la AGI}

Zhang Peng explica que en ese momento no se empeñaron en definir la AGI, porque estaba demasiado lejos y no podía definirse con claridad; en cambio, sí podía intentarse definir «cuál es el siguiente paso». Es un enfoque prudente: no sustituir un problema de la etapa actual por un concepto último. La inteligencia cognitiva le dio a Zhipu (智谱) una dirección lo bastante lejana, pero que aún podía avanzarse mediante la ingeniería.

\begin{importantbox}{El valor de los objetivos por etapas}
Una visión de largo plazo necesita definiciones por etapas. La AGI es un objetivo lejano; la inteligencia cognitiva es un objetivo de etapa que puede discutirse, que permite organizar recursos y que admite una conversión a ingeniería.
\end{importantbox}

\subsection{Resumen de la sección}

La inteligencia cognitiva fue el punto de anclaje de la dirección de Zhipu en sus inicios. Le permitió a la empresa no quedarse en la inteligencia perceptiva ni proclamar la AGI sin sustento, sino buscar una ruta hacia la siguiente generación de IA a partir de la comprensión, el razonamiento y la ingeniería del conocimiento.
